\QHeading{第1题　在$\bigcirc$和$\square$里分别填写相应的运算符号和数。}
\par\noindent


\Equation{78\times 99+78=\textcolor{answerred}{78}\times (\textcolor{answerred}{99}\textcolor{answerred}{+}\textcolor{answerred}{1})}
\par

\Equation{25\times 44=25\times\textcolor{answerred}{40}+25\times\textcolor{answerred}{4}}
\Equation{(32+15)\times 2=\textcolor{answerred}{32}\times 2+\textcolor{answerred}{15}\times 2}
\Equation{35\times 15+85\times 35=(\textcolor{answerred}{15}\textcolor{answerred}{+}\textcolor{answerred}{85})\textcolor{answerred}{\times}\textcolor{answerred}{35}}
\par\noindent

\QHeading{第2题　用简便方法计算下面各题。}

\CalcBlock[方法一：简便计算]{1}{\begin{aligned}[t]&\phantom{={}}\color{black}25\times (30+4)\\&\color{answerred}=25\times30+25\times4\\&\color{answerred}=750+100\\&\color{answerred}=850\end{aligned}}
\par
\MethodBlock{方法二：按运算顺序直接计算（对照）}{本法只作结果核对；本题要求简便计算，前面的简便解法仍保留。}{\begin{aligned}[t]&\phantom{={}}25\times (30+4)\\
&=25\times 34\\
&=850\end{aligned}}


\CalcBlock[方法一：简便计算]{2}{\begin{aligned}[t]&\phantom{={}}\color{black}126\times 67+126\times 33\\&\color{answerred}=126\times(67+33)\\&\color{answerred}=126\times100\\&\color{answerred}=12600\end{aligned}}
\MethodBlock{方法二：先处理括号或乘除，再凑整再减去多加的数}{}{\begin{aligned}[t]&\phantom{={}}126\times 67+126\times 33\\
&=8442+4158\\
&=8442+4160-2\\
&=12602-2\\
&=12600\end{aligned}}

\CalcBlock[方法一：简便计算]{3}{\begin{aligned}[t]&\phantom{={}}\color{black}63\times (2+100)\\&\color{answerred}=63\times2+63\times100\\&\color{answerred}=126+6300\\&\color{answerred}=6426\end{aligned}}
\MethodBlock{方法二：按运算顺序直接计算（对照）}{本法只作结果核对；本题要求简便计算，前面的简便解法仍保留。}{\begin{aligned}[t]&\phantom{={}}63\times (2+100)\\
&=63\times 102\\
&=6426\end{aligned}}

\CalcBlock[方法一：简便计算]{4}{\begin{aligned}[t]&\phantom{={}}\color{black}97\times 6+97\times 4\\&\color{answerred}=97\times(6+4)\\&\color{answerred}=97\times10\\&\color{answerred}=970\end{aligned}}
\MethodBlock{方法二：先处理括号或乘除，再凑整再减去多加的数}{}{\begin{aligned}[t]&\phantom{={}}97\times 6+97\times 4\\
&=582+388\\
&=582+390-2\\
&=972-2\\
&=970\end{aligned}}

\CalcBlock[方法一：简便计算]{5}{\begin{aligned}[t]&\phantom{={}}\color{black}75\times 28+25\times 28\\&\color{answerred}=(75+25)\times28\\&\color{answerred}=100\times28\\&\color{answerred}=2800\end{aligned}}
\MethodBlock{方法二：按运算顺序直接计算（对照）}{本法只作结果核对；本题要求简便计算，前面的简便解法仍保留。}{\begin{aligned}[t]&\phantom{={}}75\times 28+25\times 28\\
&=2100+700\\
&=2800\end{aligned}}

\CalcBlock[方法一：简便计算]{6}{\begin{aligned}[t]&\phantom{={}}\color{black}(125+40)\times 8\\&\color{answerred}=125\times8+40\times8\\&\color{answerred}=1000+320\\&\color{answerred}=1320\end{aligned}}
\MethodBlock{方法二：先处理括号或乘除，再凑整后补偿}{把接近整十、整百的因数凑整，再补回或减去多算的部分。}{\begin{aligned}[t]&\phantom{={}}(125+40)\times 8\\
&=165\times 8\\
&=165\times (10-2)\\
&=165\times 10-165\times 2\\
&=1650-330\\
&=1320\end{aligned}}

\par\noindent

\QHeading{第3题　解决问题。}

\Problem{① 学校买来8套演出服（每套演出服包含一件上衣和一条裤子），每件上衣125元，每条裤子80元。学校买这些演出服一共花了多少钱?}
\par


\MethodTitle{方法一：简便计算}
\Answer{\Formula{\begin{aligned}[t]&\phantom{={}}(125+80)\times 8\\&=125\times 8+80\times 8\\&=1000+640\\&=1640\text{（元）}\end{aligned}}}
\Answer{答：一共花了1640元。}

\MethodTitle{方法二：分别计算两部分，再求和}
\Answer{不先合并每份的数量，而是分别求出两部分的总数，再相加。}
\Answer{\Formula{\begin{aligned}[t]&\phantom{={}}125\times 8+80\times 8\\
&=1000+640\\
&=1640\text{（元）}\end{aligned}}}
\Answer{答：一共花了1640元。}
\Problem{② 工人师傅粉刷墙壁，上午粉刷40平方米，下午粉刷60平方米。粉刷了4天，一共粉刷了多少平方米?}

\MethodTitle{方法一}
\Answer{\Formula{\begin{aligned}[t]&\phantom{={}}(40+60)\times4\\&=100\times4\\&=400\text{（平方米）}\end{aligned}}}
\Answer{答：共粉刷400平方米。}

\MethodTitle{方法二：分别计算两部分，再求和}
\Answer{不先合并每份的数量，而是分别求出两部分的总数，再相加。}
\Answer{\Formula{\begin{aligned}[t]&\phantom{={}}40\times 4+60\times 4\\
&=160+240\\
&=400\text{（平方米）}\end{aligned}}}
\Answer{答：一共粉刷400平方米。}
